\section{Related Work}
\label{sec:related_work}

% Owner: Kuei Chen Kan (@erickan2002-del). Tracking: GitHub issue #26.
% Carries over to the final report.

\paragraph{Speculative taint tracking.}
Yu et al.~\cite{Yu:STT:2019} propose Speculative Taint
Tracking (STT), which tracks data obtained through
speculative execution and delays instructions that could
transmit that data through microarchitectural channels.
Rather than stopping all speculative execution, STT
restricts potentially leaking uses of speculative data.
Our project adopts STT as the baseline protection and
evaluates its implementation in SonicBOOM.

\paragraph{Microarchitectural improvements to STT.}
Chen et al.~\cite{Chen:uSTT:2025} identify two implementation
costs of STT: taint-propagation delay at rename and the
hardware required to delay instructions.
Their $\mu$STT design introduces an Age Matrix to simplify
taint propagation and an impede micro-op that reuses
existing data-dependency tracking to delay transmitters.
The paper evaluates the taint-propagation circuit through
RTL synthesis and the instruction-delay mechanism through
gem5 simulation on AArch64.
Our project aims to integrate both mechanisms into the
SonicBOOM RTL and evaluate their combined effects on
correctness, attack resistance, performance, and hardware
cost.

\paragraph{Secure speculation in RISC-V RTL.}
SecureBOOM~\cite{Jauch:SecureBOOM:2023} presents a
secure-by-construction methodology that combines
information-flow tracking with formal verification.
Its BOOMv3 case study demonstrates secure speculation
in a complete RTL processor.
This work shows that RTL-level secure speculation is
already feasible. Our project has a narrower focus:
implementing the specific STT and $\mu$STT mechanisms
and measuring their integration costs.
Our planned attack tests and regressions provide
experimental validation rather than the formal security
guarantees pursued by SecureBOOM.

\paragraph{Cache-focused defenses.}
SpecLFB~\cite{Cheng:SpecLFB:2024} uses a Line-Fill Buffer
and load-security checks to prevent speculative
cache side channels.
It is evaluated in gem5 and implemented in SonicBOOM,
including an FPGA prototype.
Whereas SpecLFB focuses on controlling speculative
loads and their cache effects, our project follows
taint propagation and transmitter-delay mechanisms.
SpecLFB therefore provides a related example of
evaluating a defense across simulation and RISC-V RTL.
Our contribution is the planned integration and
evaluation of $\mu$STT within SonicBOOM, rather than
a claim to introduce the first hardware Spectre defense.
